% !TEX program = tectonic
%=====================================================================
%  Criptografía y Seguridad - ITBA
%  Clase 5 — Protocolos (PKI, KDC, Needham-Schroeder, TLS, Shamir)
%  Apunte integral: teoría + práctica
%=====================================================================
\documentclass[11pt,a4paper]{article}

%---------------------------------------------------------------------
%  Paquetes
%---------------------------------------------------------------------
\usepackage{iftex}
\ifPDFTeX
  \usepackage[utf8]{inputenc}
  \usepackage[T1]{fontenc}
  \usepackage{lmodern}
\else
  \usepackage{fontspec}
\fi
\usepackage[spanish,es-tabla,es-noshorthands]{babel}
\usepackage{microtype}

\usepackage{amsmath,amssymb,amsthm}
\usepackage{mathtools}
\usepackage{geometry}
\geometry{a4paper,margin=2.4cm,headheight=22pt}

\usepackage{xcolor}
\usepackage{graphicx}
\usepackage{enumitem}
\usepackage{booktabs}
\usepackage{array}
\usepackage{tcolorbox}
\tcbuselibrary{skins,breakable}
\usepackage{fancyhdr}
\usepackage{titlesec}
\usepackage{hyperref}
\usepackage{listings}
\usepackage{caption}

\usepackage{pgfplots}
\pgfplotsset{compat=1.18}

\usepackage{tikz}
\usetikzlibrary{shapes.geometric,shapes.misc,shapes.symbols,arrows.meta,positioning,calc,
                fit,backgrounds,decorations.pathmorphing,shadows.blur,chains}

%---------------------------------------------------------------------
%  Paleta de color (inspirada en las diapositivas de la cátedra)
%---------------------------------------------------------------------
\definecolor{itbaCyan}{HTML}{6EC6D9}
\definecolor{itbaCyanDark}{HTML}{2E8B9E}
\definecolor{itbaBlue}{HTML}{AEDCEA}
\definecolor{boxBlue}{HTML}{BBD4F0}
\definecolor{slideGray}{HTML}{ECECEC}
\definecolor{practGreen}{HTML}{4F8A3D}
\definecolor{practGreenL}{HTML}{E2EFD9}
\definecolor{attackRed}{HTML}{C0392B}
\definecolor{codeBg}{HTML}{F5F5F5}

%---------------------------------------------------------------------
%  Hipervínculos
%---------------------------------------------------------------------
\hypersetup{
  colorlinks=true,
  linkcolor=itbaCyanDark,
  urlcolor=itbaCyanDark,
  citecolor=itbaCyanDark,
  pdftitle={Clase 5 - Protocolos},
  pdfauthor={Criptografia y Seguridad - ITBA}
}

%---------------------------------------------------------------------
%  Encabezados y pies
%---------------------------------------------------------------------
\pagestyle{fancy}
\fancyhf{}
\lhead{\small\textsc{Criptografía y Seguridad — ITBA}}
\rhead{\small Clase 5 · Protocolos}
\cfoot{\small\thepage}
\renewcommand{\headrulewidth}{0.4pt}
\renewcommand{\headrule}{\hbox to\headwidth{\color{itbaCyanDark}\leaders\hrule height \headrulewidth\hfill}}

%---------------------------------------------------------------------
%  Títulos de sección con barra de color (estilo diapositiva)
%---------------------------------------------------------------------
\titleformat{\section}
  {\normalfont\Large\bfseries\color{itbaCyanDark}}
  {\thesection}{0.6em}{}
  [\vspace{-0.4em}{\color{itbaCyan}\titlerule[2pt]}]
\titleformat{\subsection}
  {\normalfont\large\bfseries\color{itbaCyanDark}}
  {\thesubsection}{0.6em}{}
\titleformat{\subsubsection}
  {\normalfont\normalsize\bfseries\color{black}}
  {\thesubsubsection}{0.6em}{}

%---------------------------------------------------------------------
%  Cajas reutilizables
%---------------------------------------------------------------------
\newtcolorbox{slidebox}[1][]{
  enhanced, colback=slideGray, colframe=black!55, boxrule=0.5pt,
  arc=1pt, left=6pt,right=6pt,top=4pt,bottom=4pt, #1}

\newtcolorbox{defbox}[1]{
  enhanced, breakable, colback=itbaBlue!25, colframe=itbaCyanDark,
  boxrule=1pt, arc=2pt, left=8pt,right=8pt,top=6pt,bottom=6pt,
  fonttitle=\bfseries, coltitle=white, title={#1}}

\newtcolorbox{practbox}[1]{
  enhanced, breakable, colback=practGreenL, colframe=practGreen,
  boxrule=1pt, arc=2pt, left=8pt,right=8pt,top=6pt,bottom=6pt,
  fonttitle=\bfseries, coltitle=white, title={#1}}

\newtcolorbox{warnbox}[1]{
  enhanced, breakable, colback=attackRed!8, colframe=attackRed,
  boxrule=1pt, arc=2pt, left=8pt,right=8pt,top=6pt,bottom=6pt,
  fonttitle=\bfseries, coltitle=white, title={#1}}

%---------------------------------------------------------------------
%  Estilo de código
%---------------------------------------------------------------------
\lstdefinestyle{shell}{
  basicstyle=\ttfamily\footnotesize,
  backgroundcolor=\color{codeBg},
  keywordstyle=\color{itbaCyanDark}\bfseries,
  commentstyle=\color{practGreen}\itshape,
  stringstyle=\color{attackRed},
  breaklines=true, frame=single, rulecolor=\color{black!30},
  showstringspaces=false, columns=fullflexible, tabsize=2,
  xleftmargin=10pt, framexleftmargin=8pt}

%---------------------------------------------------------------------
%  Macros
%---------------------------------------------------------------------
\newcommand{\enc}[2]{\{#1\}\,#2}          % {M} k
\newcommand{\kdc}{\textsf{KDC}}
\newcommand{\concat}{\,\|\,}

%  Estilo común para diagramas de secuencia
\tikzset{
  lifeline/.style={thick, dotted, black!55},
  actor/.style={draw, rounded corners, fill=boxBlue, minimum width=1.3cm,
                minimum height=0.7cm, font=\bfseries},
  msg/.style={-{Stealth}, thick, itbaCyanDark},
  msgback/.style={-{Stealth}, thick, itbaCyanDark},
  atk/.style={-{Stealth}, thick, attackRed},
  note/.style={draw=black!40, fill=slideGray, rounded corners, font=\scriptsize,
               inner sep=3pt, align=center}
}

%=====================================================================
\begin{document}
%=====================================================================

%---------------------------------------------------------------------
%  Portada
%---------------------------------------------------------------------
\begin{titlepage}
\centering
\vspace*{1.2cm}

\begin{tikzpicture}
  \node[rounded corners=2pt, fill=itbaCyan, text=black, inner sep=14pt,
        minimum width=0.8\textwidth, align=center] (t)
        {\Huge\bfseries Criptografía y Seguridad};
\end{tikzpicture}

\vspace{0.6cm}
{\LARGE Criptografía: Protocolos y otros}\\[0.25cm]
{\Huge\bfseries\color{itbaCyanDark} Protocolos, PKI y Secreto Compartido}

\vspace{0.8cm}

% Ilustración: bóveda / caja fuerte (recreación del icono de portada)
\begin{tikzpicture}[scale=1.0]
  \fill[black!12] (-3.2,-2.4) rectangle (3.2,2.8);
  \draw[line width=2pt, black!55] (-3.2,-2.4) rectangle (3.2,2.8);
  \fill[itbaCyanDark!75] (-2.4,-2.0) rectangle (2.4,2.4);
  \draw[line width=3pt, black!70] (-2.4,-2.0) rectangle (2.4,2.4);
  \draw[line width=1.2pt, itbaCyan] (-2.0,-1.6) rectangle (2.0,2.0);
  \draw[line width=2pt, black!70] (0,-2.0) -- (0,2.4);
  \foreach \a in {0,45,...,315}{
    \draw[line width=2pt, black!70] (-1.1,0.2) -- ++(\a:0.6);
  }
  \draw[line width=2.5pt, black!70] (-1.1,0.2) circle (0.42);
  \fill[itbaCyan] (-1.1,0.2) circle (0.16);
  \draw[line width=2.5pt, black!70] (1.1,0.2) circle (0.42);
  \fill[itbaCyan] (1.1,0.2) circle (0.16);
  \foreach \y in {-1.2,-0.4,0.8,1.6}{
     \fill[black!70] (-2.2,\y) circle (0.07);
     \fill[black!70] (2.2,\y) circle (0.07);
  }
\end{tikzpicture}

\vfill
{\large Apunte integral — Teoría + Práctica}\\[0.2cm]
{\normalsize Clase 5 \quad·\quad Capítulo 11, \emph{Computer Security: Art and Science} (M.~Bishop) \quad·\quad RFC 8446 (TLS 1.3)}\\[0.4cm]
{\small Instituto Tecnológico de Buenos Aires (ITBA)}

\vspace{0.6cm}
{\footnotesize\itshape Documento elaborado a partir de las diapositivas oficiales de la
teórica y de la práctica, y de las transcripciones de ambas clases.}
\end{titlepage}

\tableofcontents
\newpage

%=====================================================================
\section{Introducción: del cifrado simétrico a los protocolos}
%=====================================================================

Esta clase cierra el bloque de criptografía con los \textbf{protocolos} que hacen utilizable
todo lo visto: cómo dos partes acuerdan una clave, cómo se asocia una clave pública a una
identidad (PKI), cómo se construye un canal seguro (TLS) y cómo se reparte un secreto entre
varios custodios (Shamir). La idea central es que los \emph{algoritmos} criptográficos son
piezas; los \emph{protocolos} son las reglas de intercambio de mensajes que los combinan para
lograr propiedades de seguridad frente a un adversario.

\begin{defbox}{Punto de partida (repaso)}
Un criptosistema \textbf{CCA-Secure} permite enviar información manteniendo
\textbf{confidencialidad} e \textbf{integridad}:
\[
  A \xrightarrow{\;\;\mathrm{Enc}_k(M)\;\;} B
\]
\textbf{Pero requiere que ambas partes conozcan una misma clave $k$.} De ahí surge el
problema que organiza toda la clase: \emph{¿cómo logran $A$ y $B$ compartir esa clave?}
\end{defbox}

\begin{defbox}{Intercambio de claves como protocolo $\Pi(n)$}
Un \textbf{intercambio de claves} es un protocolo $\Pi(n)$ ejecutado por dos partes:
\[
  \Pi : (n) \;\to\; (\mathrm{Tran},\, k_a,\, k_b)
\]
\begin{itemize}[leftmargin=1.4em,itemsep=1pt]
  \item No tiene entrada salvo el \textbf{parámetro de seguridad} $n$.
  \item La salida es: un conjunto de \textbf{mensajes intercambiados} ($\mathrm{Tran}$), una
        clave $k_a$ conocida solo por una parte y una clave $k_b$ conocida solo por la otra.
\end{itemize}
\textbf{Condición fundamental:} $\;k_a = k_b$ (ambas partes terminan con la misma clave).
\end{defbox}

\begin{practbox}{Aporte de la práctica — limitaciones de la clave privada (KDC)}
La práctica encuadra toda la clase en las \textbf{tres limitaciones de la criptografía de
clave privada (simétrica)}:
\begin{enumerate}[leftmargin=1.6em,itemsep=1pt]
  \item \textbf{Distribución de la clave} $\;\to\;$ se ataca con un \textbf{KDC} (Key
        Distribution Center).
  \item \textbf{Almacenamiento de la clave} $\;\to\;$ también lo alivia el KDC (cada uno
        guarda una sola clave).
  \item \textbf{Sistemas abiertos} (\emph{open systems}) $\;\to\;$ motiva la
        \emph{revolución de clave pública} (PKI).
\end{enumerate}
Con $n$ participantes que quieren hablar de a pares, un esquema todos-con-todos requiere
$\binom{n}{2}$ claves; el KDC lo reduce a $n$ claves (una por participante con el centro).
\end{practbox}

%=====================================================================
\section{Distribución simétrica de claves: el KDC}
%=====================================================================

\begin{defbox}{KDC — Key Distribution Center}
Un \textbf{punto único} en el que se concentran las relaciones de confianza. Cada entidad
$U$ comparte una clave $k_U$ \emph{solo con el KDC} ($k_A$ para $A$, $k_B$ para $B$, etc.).
Cuando $A$ quiere hablar con $B$:
\begin{enumerate}[leftmargin=1.6em,itemsep=1pt]
  \item $A \to \kdc$: ``quiero comunicarme con $B$''.
  \item El \kdc{} genera una \textbf{clave de sesión} $s$ y la entrega cifrada a cada uno:
        $\kdc \to A:\enc{s}{k_A}$ \ y \ $\kdc \to B:\enc{s}{k_B}$.
  \item $A$ y $B$ se comunican entre sí cifrando con $s$: $\enc{\dots}{s}$.
\end{enumerate}
La clave de sesión $s$ se genera \emph{ad hoc} para esa conversación.
\end{defbox}

\begin{figure}[h]
\centering
\begin{tikzpicture}[font=\small]
  \node[actor, fill=attackRed!25] (kdc) at (4,2.4) {KDC};
  \node[actor] (a) at (0,0) {A};
  \node[actor] (b) at (8,0) {B};
  \draw[msg] (a) to[bend left=18] node[above left,font=\scriptsize]{1.\ ``quiero hablar con $B$''} (kdc);
  \draw[msg] (kdc) to[bend left=18] node[below right,font=\scriptsize]{2.\ $\enc{s}{k_A}$} (a);
  \draw[msg] (kdc) to[bend left=18] node[above right,font=\scriptsize]{2.\ $\enc{s}{k_B}$} (b);
  \draw[msg] (a) -- node[below,font=\scriptsize]{3.\ usan $s$: $\enc{\dots}{s}$} (b);
\end{tikzpicture}
\caption{Esquema del KDC (recreación de la diapositiva de práctica). Ejemplo de uso real:
\textbf{Kerberos} y \textbf{Active Directory}.}
\end{figure}

\begin{center}
\renewcommand{\arraystretch}{1.25}
\begin{tabular}{>{\bfseries}p{3.2cm} p{10.2cm}}
\toprule
\textcolor{practGreen}{Ventajas} & Cada participante guarda \emph{solo} $k_U$. Hay un único sitio sensible: el KDC.\\
\textcolor{attackRed}{Desventajas} & \textbf{Centralización}: si el KDC es atacado o falla, \emph{se cae todo}. Es el único punto de confianza (y de falla).\\
\bottomrule
\end{tabular}
\end{center}

\begin{practbox}{Aclaración del profe — KDC vs.\ PKI}
El KDC y la PKI \emph{ambos} requieren confiar en un tercero, pero de naturaleza distinta:
\begin{itemize}[leftmargin=1.4em,itemsep=1pt]
  \item En el \textbf{KDC} (mundo simétrico) confío en un servicio \emph{en línea} que
        \emph{conoce todas las claves} y participa en cada conversación.
  \item En la \textbf{PKI} (mundo asimétrico) confío en una \textbf{autoridad certificante}
        que solo \emph{firma} la asociación clave--identidad; \emph{no} conoce las claves
        privadas ni participa en cada conexión.
\end{itemize}
La PKI \emph{no es aplicable} a los criptosistemas simétricos: resuelve la integridad/origen
de \emph{claves públicas}.
\end{practbox}

%=====================================================================
\section{Protocolo Needham--Schroeder (clave simétrica)}
%=====================================================================

Es el protocolo de intercambio de claves simétricas que usan \textbf{Kerberos} y
\textbf{Active Directory} (con modificaciones). Requiere un servicio centralizado (KDC) y
\textbf{genera claves de sesión entre pares}. \emph{Hipótesis}: cada entidad comparte una
clave con el KDC. A lo largo de las sucesivas ``aproximaciones'' se ven los ataques y sus
arreglos.

\medskip
\noindent\textbf{Notación.} Se escribe $e_k(M) = \enc{M}{k}$ (cifrado de $M$ bajo $k$), y
$\concat$ denota concatenación.

\subsection{Primera aproximación (ingenua)}

\begin{figure}[h]
\centering
\begin{tikzpicture}[font=\small]
  % lifelines
  \foreach \x/\n in {0/A, 8/KDC}{
    \node[actor] (\n t) at (\x,0.6) {\n};
    \node[actor] (\n b) at (\x,-3.2) {\n};
    \draw[lifeline] (\n t) -- (\n b);
  }
  \node[actor] (Bt2) at (12,0.6) {B};
  \node[actor] (Bb2) at (12,-3.2) {B};
  \draw[lifeline] (Bt2) -- (Bb2);
  % messages
  \draw[msg] (0,-0.1) -- node[above,font=\scriptsize]{1.\ $\enc{\text{Sesión }A\to B}{k_A}$} (8,-0.1);
  \draw[msgback] (8,-1.2) -- node[above,font=\scriptsize]{2.\ $\enc{s}{k_A}\concat\enc{s}{k_B}$} (0,-1.2);
  \draw[msg] (0,-2.4) -- node[above,font=\scriptsize]{3.\ $\enc{s}{k_B}$} (12,-2.4);
\end{tikzpicture}
\caption{Primera aproximación: $A$ pide la sesión, el KDC le devuelve $s$ cifrada para ambos,
$A$ reenvía a $B$ su copia. (Recreación de la diapositiva.)}
\end{figure}

\begin{warnbox}{Problemas de la primera aproximación}
\begin{itemize}[leftmargin=1.4em,itemsep=2pt]
  \item \textbf{Repetición (replay)}: un atacante graba los mensajes de $A$ hacia $B$ y luego
        reenvía $\enc{s}{k_B}$ y los siguientes. \emph{$B$ no tiene forma de saber que no
        está hablando con $A$.}
  \item \textbf{Reúso de claves (key reuse)}: un atacante graba el mensaje del KDC a $A$ y se
        lo reenvía cuando $A$ inicia otra conversación. $A$ y $B$ terminan usando la
        \emph{misma} clave de sesión vieja.
\end{itemize}
\end{warnbox}

\subsection{Segunda aproximación (con nonces)}

Se agregan \textbf{nonces} (números aleatorios de un solo uso) $r_1, r_2$ para garantizar
\textbf{frescura} (anti-replay) y autenticación mutua.

\begin{figure}[h]
\centering
\begin{tikzpicture}[font=\small]
  \node[actor] (At) at (0,1) {A};   \node[actor] (Ab) at (0,-5) {A};
  \node[actor] (Kt) at (12,1) {KDC};\node[actor] (Kb) at (12,-1.8) {KDC};
  \node[actor] (Bt) at (12,-2.6) {B}; \node[actor] (Bb) at (12,-5) {B};
  \draw[lifeline] (At) -- (Ab);
  \draw[lifeline] (Kt) -- (Kb);
  \draw[lifeline] (Bt) -- (Bb);
  \draw[msg] (0,0.4) -- node[above,font=\scriptsize]{1.\ $A\concat B\concat r_1$} (12,0.4);
  \draw[msgback] (12,-0.6) -- node[above,font=\scriptsize]{2.\ $\enc{A\concat B\concat r_1\concat s\concat \enc{A\concat s}{k_B}}{k_A}$} (0,-0.6);
  \draw[msg] (0,-2.1) -- node[above,font=\scriptsize]{3.\ $\enc{A\concat s}{k_B}$} (12,-2.1);
  \draw[msgback] (12,-3.2) -- node[above,font=\scriptsize]{4.\ $\enc{r_2}{s}$} (0,-3.2);
  \draw[msg] (0,-4.2) -- node[above,font=\scriptsize]{5.\ $\enc{r_2-1}{s}$} (12,-4.2);
\end{tikzpicture}
\caption{Segunda aproximación con nonces. (Recreación de la diapositiva.) El mensaje~2 está
cifrado con $k_A$ (clave A--KDC), el mensaje~3 contiene el ``ticket'' $\enc{A\concat s}{k_B}$
que solo $B$ abre.}
\end{figure}

\begin{defbox}{Por qué funciona — explicación mensaje a mensaje}
\begin{description}[leftmargin=2.4em,style=nextline,itemsep=2pt]
  \item[Segundo mensaje] Cifrado con la clave compartida $A$--KDC $\Rightarrow$ \emph{proviene
        del KDC}. \textbf{No es repetición} porque el $r_1$ recibido coincide con el enviado.
  \item[Cuarto mensaje] \textbf{Solo $B$ puede enviarlo} (solo $B$ obtiene $s$ del ticket).
        Avisa a $A$ de un intento de comunicación enviando $\enc{r_2}{s}$.
  \item[Quinto mensaje] $A$ confirma la comunicación devolviendo $\enc{r_2-1}{s}$. $B$ sabe
        que \emph{no es una repetición} (por $r_2$): prueba que del otro lado hay alguien que
        conoce $s$ \emph{ahora}.
\end{description}
\end{defbox}

\subsection{Ataque residual y modificación Denning--Sacco}

\begin{warnbox}{Problema: clave de sesión vieja comprometida}
Escenario: el atacante $E$ obtiene una \textbf{clave de sesión antigua} $s$. Comienza
directamente por el \textbf{tercer paso}, reenviando el ticket viejo $\enc{A\concat s}{k_B}$:
\end{warnbox}

\begin{figure}[h]
\centering
\begin{tikzpicture}[font=\small]
  \node[actor,fill=attackRed!25] (Et) at (0,0.6) {E}; \node[actor,fill=attackRed!25] (Eb) at (0,-2.8) {E};
  \node[actor] (Bt) at (8,0.6) {B}; \node[actor] (Bb) at (8,-2.8) {B};
  \draw[lifeline] (Et) -- (Eb);
  \draw[lifeline] (Bt) -- (Bb);
  \draw[atk] (0,-0.1) -- node[above,font=\scriptsize]{$\enc{A\concat s}{k_B}$ (ticket viejo)} (8,-0.1);
  \draw[msgback,attackRed] (8,-1.1) -- node[above,font=\scriptsize]{$\enc{r_2}{s}$} (0,-1.1);
  \draw[atk] (0,-2.1) -- node[above,font=\scriptsize]{$\enc{r_2-1}{s}$} (8,-2.1);
  \node[note, fill=attackRed!12, draw=attackRed] at (4,-3.4) {\textbf{¡E logra impersonar a A!}};
\end{tikzpicture}
\caption{Ataque por reúso de clave de sesión vieja. (Recreación de la diapositiva.)}
\end{figure}

\begin{defbox}{Modificación Denning--Sacco — agregar timestamps}
El arreglo propuesto es incluir un \textbf{timestamp} $T$ dentro del ticket. El ticket pasa a
ser $\enc{A\concat T\concat s}{k_B}$: $B$ verifica que $T$ sea \emph{reciente} y rechaza
tickets viejos, neutralizando el ataque anterior.
\end{defbox}

\begin{figure}[h]
\centering
\begin{tikzpicture}[font=\small]
  \node[actor] (At) at (0,1) {A};   \node[actor] (Ab) at (0,-5) {A};
  \node[actor] (Kt) at (12,1) {KDC};\node[actor] (Kb) at (12,-1.8) {KDC};
  \node[actor] (Bt) at (12,-2.6) {B}; \node[actor] (Bb) at (12,-5) {B};
  \draw[lifeline] (At) -- (Ab);
  \draw[lifeline] (Kt) -- (Kb);
  \draw[lifeline] (Bt) -- (Bb);
  \draw[msg] (0,0.4) -- node[above,font=\scriptsize]{1.\ $A\concat B\concat r_1$} (12,0.4);
  \draw[msgback] (12,-0.6) -- node[above,font=\scriptsize]{2.\ $\enc{A\concat B\concat r_1\concat s\concat \enc{A\concat \textcolor{attackRed}{T}\concat s}{k_B}}{k_A}$} (0,-0.6);
  \draw[msg] (0,-2.1) -- node[above,font=\scriptsize]{3.\ $\enc{A\concat \textcolor{attackRed}{T}\concat s}{k_B}$} (12,-2.1);
  \draw[msgback] (12,-3.2) -- node[above,font=\scriptsize]{4.\ $\enc{r_2}{s}$} (0,-3.2);
  \draw[msg] (0,-4.2) -- node[above,font=\scriptsize]{5.\ $\enc{r_2-1}{s}$} (12,-4.2);
\end{tikzpicture}
\caption{Modificación Denning--Sacco: el timestamp $T$ (en rojo) viaja dentro del ticket.
(Recreación de la diapositiva.)}
\end{figure}

\begin{practbox}{Aclaración del profe — nonces, frescura y anti-replay}
Los ataques de \textbf{replay} (repetición de mensajes) y de \textbf{key reuse} aparecen en
\emph{cualquier} protocolo (KDC, Diffie--Hellman, etc.). Las dos herramientas para garantizar
\textbf{frescura} son:
\begin{itemize}[leftmargin=1.4em,itemsep=1pt]
  \item \textbf{Nonces} ($r_1, r_2$): números aleatorios de un solo uso; quien responde con el
        nonce correcto demuestra que participa \emph{ahora}.
  \item \textbf{Timestamps} ($T$): marcas de tiempo que prueban que el mensaje es reciente y
        no una grabación vieja.
\end{itemize}
\end{practbox}

%=====================================================================
\section{El problema de las claves públicas: Man-in-the-Middle}
%=====================================================================

Con criptografía asimétrica desaparece la necesidad de compartir un secreto previo, pero
\textbf{aparece otro problema}: ¿de dónde se obtienen las claves públicas? Se confía
\emph{implícitamente} en la asociación entre una clave y la identidad de su dueño.

\begin{warnbox}{Ataques activos}
Un atacante activo puede \textbf{omitir}, \textbf{reescribir}, \textbf{reordenar} o
\textbf{repetir} mensajes. Los esquemas de intercambio ingenuos \emph{no sirven} contra
atacantes activos $\Rightarrow$ ataques \textbf{Man-in-the-Middle (MITM)}.
\end{warnbox}

\begin{figure}[h]
\centering
\begin{tikzpicture}[font=\small]
  % lo que A imagina
  \node[font=\bfseries\scriptsize] at (-0.5,1.4) {Lo que A imagina:};
  \node[actor] (a1) at (0,0.6) {A}; \node[actor] (r1) at (8,0.6) {Repo};
  \draw[msg] (a1) -- node[above,font=\scriptsize]{¿clave de $B$?} (r1);
  \draw[msgback] (8,0.2) -- node[below,font=\scriptsize]{$pk_b$} (0,0.2);
  % lo que pasa
  \node[font=\bfseries\scriptsize] at (-0.5,-1.0) {Lo que pasa en realidad:};
  \node[actor] (a2) at (0,-1.8) {A};
  \node[actor,fill=attackRed!25] (e2) at (4,-1.8) {E};
  \node[actor] (r2) at (8,-1.8) {Repo};
  \draw[msg] (a2) -- node[above,font=\scriptsize]{¿clave de $B$?} (e2);
  \draw[msg] (e2) -- node[above,font=\scriptsize]{¿clave de $B$?} (r2);
  \draw[atk] (4,-2.2) -- node[below,font=\scriptsize]{\textcolor{attackRed}{$pk_e$}} (0,-2.2);
  \draw[msgback] (8,-2.2) -- node[below,font=\scriptsize]{$pk_b$} (4,-2.2);
\end{tikzpicture}
\caption{$E$ se interpone entre $A$ y el repositorio y entrega \emph{su} clave $pk_e$ en lugar
de $pk_b$. (Recreación de las diapositivas ``Problemas'' y ``Man in the middle''.)}
\end{figure}

Una vez que $A$ tiene la clave equivocada, todo lo que cifre ``para $B$'' lo abre $E$:
\[
  \underbrace{A \xrightarrow{e_{pk_e}(M)} E}_{\text{$E$ descifra}}
  \xrightarrow{e_{pk_b}(M)} B .
\]
\textbf{El problema entra en la esfera de la administración de claves}: usar la clave
equivocada significa que \emph{no hay garantías de confidencialidad ni integridad}.

%=====================================================================
\section{Infraestructura de Clave Pública (PKI)}
%=====================================================================

\begin{defbox}{Objetivo de la PKI}
\textbf{Asociar una identidad a las claves públicas} para evitar la suplantación
(\emph{man-in-the-middle} / \emph{spoofing}). \textbf{No es aplicable a criptosistemas
simétricos.} \emph{Motivo}: la selección de claves depende de las partes; usar la clave
equivocada anula la confidencialidad y la integridad.
\end{defbox}

\subsection{Certificados}

\begin{defbox}{Certificado}
Es un mensaje que contiene \textbf{al menos}: información de identidad (p.\ ej.\ nombre),
\textbf{clave pública asociada}, fecha de emisión, intervalo de validez y \textbf{tipo de uso}
de la clave (firma de mensajes, cifrado de e-mails, firma de certificados, cifrado de sitios
web). El certificado está, a su vez, \textbf{firmado digitalmente por una autoridad
competente} (CA).
\end{defbox}

\paragraph{Uso de un certificado.} $A$ solicita el certificado de $B$ (puede proveerlo $B$ o
cualquier otra entidad). Con él, $A$ puede:
\begin{itemize}[leftmargin=1.4em,itemsep=1pt]
  \item \textbf{Verificar la identidad} comparando el \textbf{CN} (\emph{Common Name}):
        e-mails $\to$ CN${=}$dirección; servidores $\to$ CN${=}$IP o \emph{hostname};
        empresas $\to$ CN${=}$razón social.
  \item \textbf{Obtener la clave pública de $B$}: es parte del propio certificado e incluye el
        tipo de clave (p.\ ej.\ RSA-2048, DSA-EC 320).
  \item \textbf{Verificar la validez}: tipo de uso permitido y fecha de vigencia.
  \item \textbf{Verificar la integridad}: validando la \textbf{firma digital} de la autoridad.
\end{itemize}

\begin{slidebox}
\textbf{El gran ``pero''}: la validación de una firma digital \emph{requiere una clave
pública}. ¿Cómo se obtiene la clave pública \emph{de la autoridad}? $\Rightarrow$ cadenas de
firmas.
\end{slidebox}

\subsection{Cadenas de firmas y CA raíces}

\begin{defbox}{Cadenas de certificados}
\textbf{Las autoridades certificantes tienen, a su vez, un certificado.} Se acostumbra incluir
el certificado de la autoridad junto con cada certificado. ¿Quién valida el certificado de la
CA? \textbf{Otra CA}.
\begin{itemize}[leftmargin=1.4em,itemsep=1pt]
  \item \textbf{CA raíces}: firman su \emph{propio} certificado (auto-firmado). Son el
        \textbf{punto de confianza} del sistema.
  \item \textbf{Concepto de raíz}: confiar en una única CA que \emph{delega} en otras la
        capacidad de firmar certificados.
\end{itemize}
\textbf{Ejemplos} (la cadena se concatena $\concat$ del sujeto hacia la raíz):
\[
  \text{Cert. de }A\text{ firmado por CA3:}\quad C_a = C'_a \concat C_{\text{CA3}} \concat C_{\text{CA2}} \concat C_{\text{CA1}}
\]
\[
  \text{Cert. de }B\text{ firmado por CA1:}\quad C_b = C'_b \concat C_{\text{CA1}}
\]
\end{defbox}

\begin{figure}[h]
\centering
\begin{tikzpicture}[font=\small, node distance=6mm]
  \node[actor, fill=itbaCyan] (root) {CA1 (raíz, auto-firmada)};
  \node[actor, below=of root] (ca2) {CA2};
  \node[actor, below=of ca2] (ca3) {CA3};
  \node[actor, below=of ca3] (a) {Cert.\ de A};
  \draw[msg] (root) -- node[right,font=\scriptsize]{firma} (ca2);
  \draw[msg] (ca2) -- node[right,font=\scriptsize]{firma} (ca3);
  \draw[msg] (ca3) -- node[right,font=\scriptsize]{firma} (a);
\end{tikzpicture}
\caption{Cadena de confianza: la raíz delega y firma hacia abajo hasta el certificado final.}
\end{figure}

\begin{defbox}{Validación entre CA distintas}
\textbf{No existe una única CA raíz.} Si $A$ y $B$ tienen CA diferentes, ¿cómo valida $A$ el
certificado de $B$?
\begin{itemize}[leftmargin=1.4em,itemsep=1pt]
  \item \textbf{Opción 1}: $A$ confía en la CA de $B$.
  \item \textbf{Opción 2}: las CA \textbf{se certifican entre sí} (cada una emite un
        certificado de la identidad de la otra).
\end{itemize}
\textbf{En la práctica}: existe una \textbf{lista de CA reconocidas} embebida en el sistema
operativo, en los navegadores y en \emph{runtimes} como la JVM.
\end{defbox}

\subsection{Certificados X.509}

\begin{defbox}{X.509 — estándar de certificados digitales}
Incluye: \textbf{versión}, número de serie, identificador del algoritmo de firma, nombre del
emisor (CA), intervalo de validez, nombre del sujeto (CN), \textbf{clave pública del sujeto} y
\textbf{firma} (firma del hash de todo lo anterior). Forman cadenas de certificados.
\end{defbox}

\begin{slidebox}
\textbf{Ejemplo (recorte de un X.509 real, formato OpenSSL):}
\begin{lstlisting}[style=shell]
Version: 3 (0x2)
Serial Number: 1 (0x1)
Signature Algorithm: md5WithRSAEncryption
Issuer:  C=ZA, ST=Western Cape, ... CN=Thawte Server CA/Email=...
Validity
    Not Before: Aug  1 00:00:00 1996 GMT
    Not After : Dec 31 23:59:59 2020 GMT
Subject: C=ZA, ... CN=Thawte Server CA/Email=...
Subject Public Key Info:
    Public Key Algorithm: rsaEncryption
    RSA Public Key: (1024 bit) ... Exponent: 65537 (0x10001)
X509v3 extensions:
    X509v3 Basic Constraints: critical  ->  CA:TRUE
Signature Algorithm: md5WithRSAEncryption  07:fa:4c:69:...
\end{lstlisting}
\end{slidebox}

\begin{defbox}{Verificación de un certificado X.509}
\begin{enumerate}[leftmargin=1.6em,itemsep=1pt]
  \item \textbf{Obtener la clave pública del emisor}: de la cadena de certificados, o del
        sistema si es raíz. Si hace falta, verificar \emph{recursivamente} el certificado del
        emisor.
  \item \textbf{Verificar la integridad}: con el algoritmo especificado y la clave pública del
        emisor (validar la firma).
  \item \textbf{Verificar el intervalo de validez}: debe estar vigente; la CA debe haber
        estado vigente al \emph{comienzo} del período de vigencia.
  \item \textbf{Verificar la identidad}: comparar el CN con el que espera la aplicación (según
        el uso del certificado).
  \item \textbf{Verificar el uso}: que el certificado esté autorizado para el uso pretendido.
\end{enumerate}
\end{defbox}

\begin{practbox}{Ejercicio (práctica/teoría) — cadena \texttt{www.itba.edu.ar}}
La diapositiva muestra una cadena real de 4 certificados. \textbf{Las validaciones necesarias}
para aceptar el del servidor:
\begin{enumerate}[leftmargin=1.6em,itemsep=1pt]
  \item \textbf{Servidor} (CN \texttt{www.itba.edu.ar}, SAN \texttt{www.itba.edu.ar}, emisor
        \emph{Amazon}): chequear que el CN/SAN coincide con el host pretendido; vigencia
        (6/6/2017–7/7/2018); validar su firma con la clave pública de \emph{Amazon}; uso
        permitido para TLS de servidor.
  \item \textbf{Intermedio} (CN \emph{Amazon}, emisor \emph{Amazon Root CA 1}): validar su
        firma con la clave de \emph{Amazon Root CA 1}; vigencia; que tenga atributo de CA.
  \item \textbf{Intermedio/raíz} (CN \emph{Amazon Root CA 1}, emisor \emph{Starfield Services
        Root CA - G2}): validar su firma con la clave de Starfield; vigencia.
  \item \textbf{Raíz} (CN \emph{Starfield Services Root CA - G2}, \textbf{auto-firmada}): debe
        estar en la \textbf{lista de CA confiables} del SO/navegador; vigencia (1/9/2009–28/6/2034).
\end{enumerate}
En cada eslabón: la clave pública del \emph{emisor} valida la firma del \emph{sujeto}, hasta
llegar a una raíz confiada. Además, conviene chequear que ninguno esté \textbf{revocado}.
\end{practbox}

\subsection{Revocación: CRL y OCSP}

\begin{defbox}{Revocación de claves}
Necesidad de invalidar una clave \textbf{antes} de su expiración: fue averiguada por un
atacante, o hay un cambio anticipado (cambio de dueño). \emph{Problemas}: no revocar claves
que no deben revocarse (evitar revocaciones sin autorización) y \textbf{propagar} la
información para evitar futuras comunicaciones con esa clave.
\end{defbox}

\begin{itemize}[leftmargin=1.4em,itemsep=2pt]
  \item \textbf{CRL (Listas de revocación)}: listas de certificados revocados, análogas a las
        listas de tarjetas de crédito robadas. \emph{Actual}: vigentes y revocados;
        \emph{histórico}: expirados y revocados. En X.509 \textbf{solo el emisor de un
        certificado puede revocarlo}; se agrega al CRL de la CA y puede validarse
        \emph{offline} u \emph{online}.
  \item \textbf{OCSP (Online Certificate Status Protocol)}: certifica que una firma está
        ``activa''. Servicio en línea con API estandarizada, administrado por las CA; el
        resultado está \textbf{firmado} (es un certificado de \emph{corta duración}: 1~hora a
        7~días). \textbf{OCSP Stapling}: el resultado se incluye en la cadena, de modo que el
        cliente no consulta una fuente externa (recomendado para sitios de alto tráfico).
\end{itemize}

\begin{slidebox}
\textbf{Resumen de certificados:} asocian una clave pública con una identidad; se pueden
validar \emph{offline} (¡cuidado con las revocaciones!); están certificados por otra entidad;
forman cadenas de confianza; \textbf{resuelven la integridad de las claves públicas}.
\end{slidebox}

%=====================================================================
\section{Canal seguro: SSL / TLS}
%=====================================================================

\begin{defbox}{Canal seguro}
\textbf{Confidencialidad e integridad sobre un canal inseguro.} SSL/TLS ofrece seguridad a
\textbf{nivel transporte}, sobre un canal confiable (TCP). Provee \textbf{confidencialidad},
\textbf{integridad} y \textbf{autenticación} de los participantes (origen y destino).
\begin{itemize}[leftmargin=1.4em,itemsep=1pt]
  \item \textbf{SSL} (Secure Socket Layer): versión inicial, creada por Netscape.
  \item \textbf{TLS} (Transport Layer Security): evolución de SSL (TLS 1.2 $\equiv$ SSL 3.3),
        mejora deficiencias; es el \textbf{estándar actual}. Depende de la \textbf{PKI}.
\end{itemize}
\end{defbox}

\begin{figure}[h]
\centering
\begin{tikzpicture}[font=\small, node distance=4mm]
  % insegura
  \node[draw,fill=slideGray] (a1) {Aplicación};
  \node[draw,fill=slideGray, below=8mm of a1] (t1) {Transporte};
  \node[draw,fill=slideGray, right=5cm of a1] (a2) {Aplicación};
  \node[draw,fill=slideGray, below=8mm of a2] (t2) {Transporte};
  \draw[msg] (a1)--(t1); \draw[msg] (t2)--(a2);
  \draw[dotted,thick] (t1) -- node[above,font=\scriptsize]{$m$ (en claro)} (t2);
  \node[font=\bfseries\scriptsize] at ($(a1)!0.5!(a2)+(0,0.9)$) {Conexión normal, insegura};
\end{tikzpicture}

\vspace{0.4cm}

\begin{tikzpicture}[font=\small, node distance=2mm]
  \node[draw,fill=slideGray] (a1) {Aplicación};
  \node[draw,fill=practGreenL, below=7mm of a1] (s1) {SSL};
  \node[draw,fill=slideGray, below=0mm of s1] (t1) {Transporte};
  \node[draw,fill=slideGray, right=5cm of a1] (a2) {Aplicación};
  \node[draw,fill=practGreenL, below=7mm of a2] (s2) {SSL};
  \node[draw,fill=slideGray, below=0mm of s2] (t2) {Transporte};
  \draw[msg] (a1)--(s1); \draw[msg] (s2)--(a2);
  \draw[dotted,thick] (t1) -- node[above,font=\scriptsize]{\texttt{(\$!)\$J"=?*} (cifrado)} (t2);
  \node[font=\bfseries\scriptsize] at ($(a1)!0.5!(a2)+(0,0.9)$) {Conexión con SSL, protegida};
\end{tikzpicture}
\caption{SSL se inserta entre aplicación y transporte. (Recreación de la diapositiva
``Concepto''.)}
\end{figure}

\subsection{TLS Record}
TLS recibe un mensaje a enviar, lo \textbf{divide en bloques} de no más de $2^{16}$ bytes,
\textbf{comprime} cada bloque y \textbf{calcula su hash}, \textbf{cifra} cada bloque y su hash,
y lo \textbf{envía} al servicio inferior (TCP). Cada \emph{record} lleva un \textbf{header}.

\subsection{Tipos de mensajes y material criptográfico}

\begin{center}
\renewcommand{\arraystretch}{1.2}
\begin{tabular}{>{\bfseries}p{3.8cm} p{9.6cm}}
\toprule
TLS handshake & Configuración inicial, autenticación y negociación del material criptográfico.\\
TLS application data & Información del protocolo encapsulado.\\
TLS alert & Informa eventos y problemas (advertencia o fatal).\\
TLS change cipher spec & Concluye una renegociación de claves.\\
\bottomrule
\end{tabular}
\end{center}

\begin{slidebox}
\textbf{Criptografía soportada (lo vigente en azul en las slides):}
\begin{itemize}[leftmargin=1.4em,itemsep=0pt]
  \item \emph{Intercambio de clave}: \textbf{RSA}, \textbf{DH con certificado RSA o DSS}, DH
        anónimo, \textbf{ECDH} (DH sobre curvas elípticas), Kerberos. (Variantes obsoletas de
        512 bits por restricciones de exportación de EE.UU.)
  \item \emph{Cifrado simétrico}: ChaCha20, 3DES-EDE-CBC, \textbf{AES} (CBC, CCM,
        \textbf{GCM}), IDEA-CBC, ARIA. (Obsoletos: RC4, DES-CBC, claves de 40 bits.)
  \item \emph{Hash/MAC}: HMAC-MD5, HMAC-SHA1, \textbf{HMAC-SHA2 (256/384)}, \textbf{AEAD}
        (Authenticated Encryption with Associated Data).
  \item \emph{Autenticación}: firmas RSA, \textbf{firmas DSS}.
\end{itemize}
\end{slidebox}

\begin{defbox}{Sesión vs.\ conexión TLS}
\textbf{Sesión}: asociación entre dos pares; puede soportar múltiples conexiones. Contiene:
identificador único, certificado X.509v3 del otro extremo (opcional), método de compresión,
método de cifrado y MAC acordados, y el \textbf{Master Secret} (clave de sesión compartida de
48 bytes).\\[3pt]
\textbf{Conexión}: describe cómo intercambiar datos. Contiene: secuencia de bits aleatoria de
calidad criptográfica, claves de escritura (distintas) para cada lado, claves para los MAC,
IVs si hacen falta, y números de secuencia para cliente y servidor.
\end{defbox}

\subsection{TLS Handshake (diagrama de secuencia)}

\begin{figure}[h]
\centering
\begin{tikzpicture}[font=\small]
  \node[actor] (Ct) at (0,1) {C}; \node[actor] (Cb) at (0,-9.5) {C};
  \node[actor] (St) at (11,1) {S}; \node[actor] (Sb) at (11,-9.5) {S};
  \draw[lifeline] (Ct) -- (Cb);
  \draw[lifeline] (St) -- (Sb);
  % parte 1
  \draw[msg] (0,0.3) -- node[above,font=\scriptsize]{\textbf{ClientHello}: $V_c\concat r_1\concat S_{id}\concat \text{Ciphers}\concat \text{Comps}$} (11,0.3);
  \draw[msgback] (11,-0.5) -- node[above,font=\scriptsize]{\textbf{ServerHello}: $V\concat r_2\concat S_{id}\concat \text{Cipher}\concat \text{Comp}$} (0,-0.5);
  % parte 2
  \draw[msgback] (11,-1.5) -- node[above,font=\scriptsize]{\textbf{Certificate} (si el servidor se autentica)} (0,-1.5);
  \draw[msgback] (11,-2.3) -- node[above,font=\scriptsize]{\textbf{ServerKeyExchange}: $p\concat \enc{\mathrm{hash}(r_1\concat r_2\concat p)}{k_s}$} (0,-2.3);
  \draw[msgback] (11,-3.1) -- node[above,font=\scriptsize]{\textbf{CertificateRequest} (si pide cert.\ al cliente): $\text{ctype}\concat \text{CAs}$} (0,-3.1);
  \draw[msgback] (11,-3.9) -- node[above,font=\scriptsize]{\textbf{ServerHelloDone}: \{\,\}} (0,-3.9);
  % parte 3
  \draw[msg] (0,-4.9) -- node[above,font=\scriptsize]{\textbf{ClientCertificate} (si fue solicitado)} (11,-4.9);
  \draw[msg] (0,-5.7) -- node[above,font=\scriptsize]{\textbf{ClientKeyExchange} (RSA): $\enc{V\concat \mathrm{PRE}_{rand}}{k_s}$ \ \ ó (DH): $\enc{g^b \bmod p}{k_s}$} (11,-5.7);
  % parte 4
  \draw[msg] (0,-6.7) -- node[above,font=\scriptsize]{\textbf{ChangeCipherSpec}: \{\,\}} (11,-6.7);
  \node[note] at (5.5,-7.15) {A partir de aquí el cliente usa los parámetros negociados};
  \draw[msg] (0,-7.9) -- node[above,font=\scriptsize]{\textbf{Finished}: $\enc{h(\text{master}\concat \text{opad}\concat h(\text{msgs}\concat \text{master}\concat \text{ipad}))}{}$} (11,-7.9);
  \draw[msgback] (11,-8.7) -- node[above,font=\scriptsize]{\textbf{ChangeCipherSpec} + \textbf{Finished} (servidor)} (0,-8.7);
  \node[note] at (5.5,-9.15) {A partir de aquí el servidor usa los parámetros negociados};
\end{tikzpicture}
\caption{Handshake completo de TLS. (Recreación de las 4 partes de las diapositivas.)}
\end{figure}

\begin{defbox}{Parámetros del handshake}
\begin{itemize}[leftmargin=1.4em,itemsep=1pt]
  \item $V_c$: versión del cliente; \ $V = \min(V_c, V_{\text{servidor}})$.
  \item $r_1, r_2$: \textbf{nonces} (timestamp + 28 bytes aleatorios); \textbf{evitan replay}.
  \item $S_{id}$: ID de sesión ($0$ para iniciar una nueva).
  \item Ciphers/Cipher, Comps/Comp: \emph{stacks} criptográficos y métodos de compresión
        ofrecidos vs.\ el elegido.
  \item $p$: parámetros criptográficos ($e,n$ para RSA; $p,g,g^a$ para DH).
  \item $k_s$: clave privada del servidor (correspondiente a la pública del certificado).
  \item \textbf{Anti-downgrade}: el cliente reenvía la versión $V$ originalmente informada
        dentro del ClientKeyExchange para \textbf{prevenir downgrade attacks}.
  \item DH: $\mathrm{PRE} = g^{ab} \bmod p$ es el \emph{pre-master secret}.
\end{itemize}
El \textbf{Master Secret} se deriva combinando $\mathrm{pre}$, $r_1$, $r_2$ con MD5/SHA:
\[
  \text{Master} = \mathrm{MD5}(\mathrm{pre}\concat \mathrm{SHA}('A'\concat \mathrm{pre}\concat r_1\concat r_2))
  \concat \mathrm{MD5}(\mathrm{pre}\concat \mathrm{SHA}('BB'\concat \cdots)) \concat \cdots
\]
\end{defbox}

\begin{itemize}[leftmargin=1.4em,itemsep=2pt]
  \item \textbf{Change Cipher Spec}: usado durante el handshake, puede aparecer en cualquier
        momento; implica \textbf{renegociación} de las claves de sesión; lo solicita cliente o
        servidor; \emph{no tiene contenido} (es un tipo de mensaje).
  \item \textbf{Alert}: eventos fuera de banda, de \emph{advertencia} o \emph{fatales} (un
        fatal invalida la conexión). \texttt{CloseNotify} indica el fin de una sesión. Errores
        fatales: mensaje no esperado, MAC incorrecto, error al descomprimir, error de
        handshake, parámetro ilegal. Advertencias: no hay certificado / no válido / no
        soportado / expirado / revocado.
\end{itemize}

\begin{slidebox}
\textbf{Panorama:} TLS es el estándar de comunicación segura en servicios web (HTTPS); depende
de la PKI; no suele usarse para validar clientes; soportado por todos los navegadores (con
distinta respuesta frente a alertas). Problemas de diseño en los clientes le hicieron perder
parte de su utilidad, pero \emph{sigue siendo la alternativa más adoptada}.
\end{slidebox}

%=====================================================================
\section{Criptografía de umbrales (Threshold Cryptography)}
%=====================================================================

\begin{defbox}{Criptosistema de umbrales}
Es una terna de algoritmos:
\[
\begin{aligned}
  \text{Gen}: ()&\to K^N &&\text{(generador de clave)}\\
  \text{Enc}: K^T \times P &\to C &&\text{(cifrado)}\\
  \text{Dec}: K^T \times C &\to P &&\text{(descifrado)}
\end{aligned}
\]
\textbf{Propiedad básica:} para todo $m$ y $k^n$ válidos, para todo
$k' = \{k^n\}_t$ y $k'' = \{k^n\}_t$ se cumple $\;d_{k''}(e_{k'}(m)) = m$.\\[4pt]
\textbf{Conceptos nuevos:}
\begin{itemize}[leftmargin=1.4em,itemsep=0pt]
  \item $N$ = \textbf{sombras} o cantidad de claves disponibles.
  \item $T$ = \textbf{umbral}: cantidad de claves \emph{necesarias}.
  \item $K^X$: $X$ claves, cada una perteneciente a $K$; \ $\{K^x\}_n$: subconjunto de $n$
        claves cualesquiera de $K^X$.
\end{itemize}
\end{defbox}

La idea: \emph{cualquier} subconjunto de $T$ de las $N$ sombras alcanza para operar, pero con
menos de $T$ no se puede.

%=====================================================================
\section{Secreto compartido de Shamir}
%=====================================================================

\begin{defbox}{Método de Shamir — esquema $(t,n)$ threshold}
\textbf{Principio:} un polinomio de grado $t-1$ queda \textbf{unívocamente determinado por su
evaluación en $t$ puntos distintos}. El \textbf{secreto es el polinomio} (en particular,
$s = P(0)$, el término independiente) y los \textbf{puntos son las sombras}.\\[4pt]
\textbf{Construcción (reparto):} sea $s$ el secreto a compartir.
\begin{itemize}[leftmargin=1.4em,itemsep=1pt]
  \item Se fija $a_0 = s$ y se eligen $a_1,\dots,a_{t-1} \leftarrow \{0,\dots,p-1\}$ al azar.
  \item Se arma $P(x) = a_{t-1}x^{t-1} + \cdots + a_1 x + a_0 \pmod p$, con $p > s$ y $p > n$
        ($p$ primo).
  \item Las sombras son $P(1), P(2), \dots, P(n)$ (se pueden generar muchas: $n$ a elección).
\end{itemize}
\textbf{Reconstrucción:} con $t$ sombras cualesquiera $(s_{i_1},\dots,s_{i_t})$, $s_{i_a}=P(i_a)$,
se interpola por \textbf{Lagrange}:
\[
  P(x) = \sum_{a=1}^{t} \left( s_{i_a} \prod_{\substack{b=1\\ b\neq a}}^{t}
          \frac{x - i_b}{\,i_a - i_b\,} \right)
\]
y se evalúa $\;s = P(0)$.
\end{defbox}

\subsection{Intuición geométrica}

\begin{practbox}{Aporte de la práctica — la recta, la parábola y las imágenes}
La práctica desarrolla la intuición:
\begin{itemize}[leftmargin=1.4em,itemsep=1pt]
  \item \textbf{Grado 1 (recta) $\Rightarrow$ esquema $(2,n)$}: hacen falta \emph{2} puntos
        para trazar una recta. Con un solo punto pasan \textbf{infinitas} rectas $\Rightarrow$
        secreto indeterminado.
  \item \textbf{Grado 2 (parábola) $\Rightarrow$ esquema $(3,n)$}: hacen falta \emph{3}
        puntos. Con 1 o 2 puntos, hay miles de parábolas posibles.
  \item En general: \textbf{número de sombras necesarias $=$ grado del polinomio $+\,1 = t$}.
        Se pueden repartir \emph{muchas} sombras (a todo un curso), pero juntando cualesquiera
        $t$ se recupera el secreto, y con menos de $t$ \emph{no se sabe nada}.
  \item \textbf{Secreto compartido por imágenes}: una imagen secreta se divide en ``capas''
        que se reparten como imágenes; superponiendo $t$ de ellas se reconstruye la imagen
        original (misma idea $(t,n)$ aplicada a píxeles).
\end{itemize}
\end{practbox}

\begin{center}
\begin{tikzpicture}
\begin{axis}[
  width=12cm, height=7cm,
  scale only axis,
  xlabel={$x$}, ylabel={$P(x)$},
  axis lines=middle,
  xmin=-0.6, xmax=4.6, ymin=-0.5, ymax=10.5,
  xtick={0,1,2,3,4}, ytick={0,2,4,6,8,10},
  enlargelimits=false, clip=true,
  legend pos=north west, legend cell align=left,
  every axis legend/.append style={font=\footnotesize},
  label style={font=\small}, tick label style={font=\footnotesize}
]
  % polinomio de grado t-1=2 (curva conceptual, valores en rango)
  \addplot[domain=-0.4:4.4, samples=120, thick, itbaCyanDark] {0.5*x*x-0.5*x+7};
  \addlegendentry{$P(x)$, grado $t{-}1=2$ (pasa por las sombras)}
  % curva "fantasma": otra parabola por (0,7) y otro punto -> ilustra ambiguedad con < t puntos
  \addplot[domain=-0.4:4.4, samples=120, dashed, attackRed] {-0.6*x*x+1.5*x+7};
  \addlegendentry{otra curva posible con menos de $t$ puntos}
  % sombras P(1)=7, P(2)=8, P(3)=10 (sobre la curva azul)
  \addplot[only marks, mark=*, mark size=2.5pt, black]
    coordinates {(1,7) (2,8) (3,10)};
  \node[anchor=north west, font=\scriptsize] at (axis cs:1.05,6.9) {sombras};
  % secreto
  \addplot[only marks, mark=square*, mark size=3.2pt, practGreen] coordinates {(0,7)};
  \node[practGreen, anchor=south west, font=\bfseries\footnotesize] at (axis cs:0.08,7.1) {$S=P(0)$};
\end{axis}
\end{tikzpicture}
\captionof{figure}{Secreto compartido de Shamir (esquema conceptual): el polinomio de grado
$t-1$ pasa por las sombras (puntos negros); el secreto es $S=P(0)$ (cuadrado verde). Con
\emph{menos} de $t$ puntos hay \textbf{infinitas} curvas posibles (la curva roja punteada
ilustra una alternativa que pasa por el mismo $S$).}
\end{center}

\subsection{Ejemplo de cifrado (reparto) — esquema $(3,5)$}

\begin{practbox}{Ejemplo resuelto (diapositiva ``Ejemplo Cifrado'')}
\textbf{Secreto} $s=7$, esquema $(3,5)$, en $\mathbb{Z}_{11}$. Se elige el polinomio de grado
$t-1=2$:
\[
  P(x) = 5x^2 + 3x + 7 \pmod{11}
\]
Se calculan las sombras $P(1),\dots,P(5)$:
\[
\begin{aligned}
  P(1) &= 5+3+7 = 15 \equiv 4 \pmod{11}\\
  P(2) &= 20+6+7 = 33 \equiv 0 \pmod{11}\\
  P(3) &= 45+9+7 = 61 \equiv 6 \pmod{11}\\
  P(4) &= 80+12+7 = 99 \equiv \mathbf{0} \pmod{11}\\
  P(5) &= 125+15+7 = 147 \equiv 4 \pmod{11}
\end{aligned}
\]
\textbf{Sombras:} $(1,4),\ (2,0),\ (3,6),\ (4,0),\ (5,4)$.
\end{practbox}

\begin{warnbox}{Errata de la diapositiva — $P(4)$}
La diapositiva original escribe $P(4)=2$, pero $80+12+7=99=9\cdot 11$, luego
$P(4)\equiv \mathbf{0}\pmod{11}$, \emph{no} $2$. Es un error aritmético de la slide. La
reconstrucción de la slide usa $(2,0),(3,6),(5,4)$ —que \emph{no} incluyen la sombra $4$—, por
lo que el resultado $s=7$ sigue siendo correcto. (Verificado aparte: las sombras correctas son
$4,0,6,0,4$.)
\end{warnbox}

\subsection{Ejemplo de descifrado (reconstrucción por Lagrange)}

\begin{practbox}{Ejemplo resuelto (diapositiva ``Ejemplo Descifrado'')}
Tomando tres sombras cualesquiera, p.\ ej.\ $(2,0),(3,6),(5,4)$, se interpola por Lagrange
(todo $\bmod\ 11$):
\[
  P(x) = \sum_{a=1}^{t} \left( s_{i_a} \prod_{\substack{b=1\\ b\neq a}}^{t}
          \frac{x - i_b}{i_a - i_b} \right)
\]
\[
  P(x) = \left[\, 0\cdot\frac{(x-3)(x-5)}{(2-3)(2-5)}
              + 6\cdot\frac{(x-2)(x-5)}{(3-2)(3-5)}
              + 4\cdot\frac{(x-2)(x-3)}{(5-2)(5-3)} \,\right]
\]
\[
  P(x) = \left[\, 0 \;-\; 3\,(x^2-7x+10) \;+\; (4/6)(x^2-5x+6) \,\right] \bmod 11
\]
\[
  P(x) = 5x^2 + 3x + 7 \qquad\Longrightarrow\qquad \boxed{\,s = P(0) = 7\,}
\]
Recordando que en $\mathbb{Z}_{11}$ la división es multiplicación por el inverso: el término
medio $\frac{6}{(1)(-2)} = \frac{6}{-2}$ y $4/6$ se resuelven con la tabla de inversos
modulares.
\end{practbox}

\begin{practbox}{Aporte de la práctica — aritmética modular e implementación}
La práctica desarrolla un \emph{segundo} ejemplo equivalente: $s=7$, esquema $(3,n)$ con
$P(x)=5x^2+2x+7$ (variante mostrada en clase) en $\mathbb{Z}_{11}$, resuelto de dos maneras:
\begin{itemize}[leftmargin=1.4em,itemsep=1pt]
  \item \textbf{Por sistema de ecuaciones (Gauss)}: plantear $a_0 + a_1 x + a_2 x^2$ para cada
        sombra y resolver el sistema $3\times 3$ con eliminación gaussiana \emph{módulo 11}
        (pivoteo, restando filas). Cómodo en papel, incómodo en código (arrastra las $x$).
  \item \textbf{Por interpolación de Lagrange}: más fácil de \emph{implementar} (no arrastra
        las $x$). Es la recomendada para el TP.
\end{itemize}
\textbf{Claves de la aritmética modular} (todo $\bmod\ 11$):
\begin{itemize}[leftmargin=1.4em,itemsep=1pt]
  \item \emph{Dividir} por $d$ $=$ \emph{multiplicar} por $d^{-1}$. Ej.: $10^{-1}=10$ (pues
        $10\cdot 10 = 100 \equiv 1$); $5^{-1}=9$ (pues $5\cdot 9 = 45 \equiv 1$);
        $8^{-1}=7$ (pues $8\cdot 7 = 56 \equiv 1$).
  \item Conviene tener a mano la \textbf{tabla de inversos modulares}.
\end{itemize}
\textbf{En el TP} se usa $\mathbb{Z}_{251}$ (el primo más grande que entra en un byte) en
lugar de $\mathbb{Z}_{11}$.
\end{practbox}

%=====================================================================
\section{Cross-check teoría vs.\ práctica}
%=====================================================================

\begin{center}
\renewcommand{\arraystretch}{1.3}
\begin{tabular}{>{\bfseries}p{3.6cm} p{4.7cm} p{4.7cm}}
\toprule
Tema & Teoría (slides) & Práctica (slides + transcript) \\
\midrule
KDC & Mencionado como base de Needham-Schroeder & \textcolor{practGreen}{Desarrolla el flujo $A\to$KDC$\to A,B$, pros/contras, Kerberos}\\
Needham-Schroeder & Las 3 aproximaciones + Denning-Sacco, paso a paso & Encuadre como limitación de clave privada\\
MITM / Masquerading & MITM sobre claves públicas & \textcolor{practGreen}{Distingue MITM (activo, ambos extremos) de Masquerading (un solo extremo); replay y key reuse}\\
PKI / Certificados & X.509, cadenas, CRL/OCSP, verificación & \textcolor{practGreen}{Certificado vs.\ clave pública ``pelada''; PKI Argentina}\\
TLS & Handshake completo, record, alert & \textcolor{practGreen}{Resumen conf./aut./integridad; DH \emph{autenticado}; RFC 2246/8446}\\
Shamir & Construcción + Lagrange + ejemplo $(3,5)$ & \textcolor{practGreen}{Intuición recta/parábola, imágenes, Gauss vs.\ Lagrange, inversos modulares, $\mathbb{Z}_{251}$ en TP}\\
\bottomrule
\end{tabular}
\end{center}

\noindent\textbf{Lo que agrega la práctica:} (1) el encuadre KDC como solución a la
distribución/almacenamiento de claves; (2) el detalle del flujo del KDC con pros y contras;
(3) la distinción precisa MITM vs.\ Masquerading y el rol de replay/key-reuse; (4) la
intuición geométrica de Shamir (recta/parábola/imágenes); (5) la resolución \emph{completa}
del ejemplo de reconstrucción y la mecánica de la aritmética modular (inversos) con ambos
métodos (Gauss y Lagrange).

%=====================================================================
\section*{Lectura recomendada}
\addcontentsline{toc}{section}{Lectura recomendada}
%=====================================================================
\begin{center}
\begin{tcolorbox}[enhanced,colback=itbaBlue!20,colframe=itbaCyanDark,boxrule=1pt,
  arc=3pt,width=0.9\textwidth,halign=center]
\large\textbf{Capítulo 11} — \emph{Computer Security: Art and Science} (Matt Bishop)\\[6pt]
\normalsize \textbf{RFC 8446} — TLS v1.3 \quad(\url{http://tools.ietf.org/html/rfc8446})\\[3pt]
\textbf{RFC 2246} — TLS (primera versión)\\[3pt]
\emph{Applied Cryptography} — Bruce Schneier\\[3pt]
Vulnerabilidad en la renegociación de claves: \url{http://www.g-sec.lu/practicaltls.pdf}
\end{tcolorbox}
\end{center}

%---------------------------------------------------------------------
\vfill
\begin{center}\small\itshape\color{black!60}
Apunte integral de la Clase 5 — Protocolos · Criptografía y Seguridad, ITBA.\\
Síntesis de teoría (diapositivas + transcripción) y práctica (KDC, Needham-Schroeder,
Denning-Sacco, TLS y Secreto Compartido de Shamir).
\end{center}

\end{document}
